\begin{table}[!htbp]
\centering
\small
\caption{Distances $|x_{n,m}-m|$ for $m=0,1,2$ in the main experiment ($\alpha=-5/2$, $j=10$, $\lambda=100$). Theorem~\ref{thm:rates}(2) gives the bound $O(q^{n})$ for every $q\in(1-c_\mu,1)$. Two significant digits are shown; the values were stable to $8$ significant digits under doubling of the precision, and each is certified to relative accuracy $10^{-6}$ in exact arithmetic.}
\label{tab:integer}
\renewcommand{\arraystretch}{1.15}
\begin{tabular}{rcccccc}
\toprule
& \multicolumn{3}{c}{Charlier, $\mu=5$} & \multicolumn{3}{c}{Meixner, $\gamma=2$, $\mu=1/2$}\\
\cmidrule(lr){2-4}\cmidrule(lr){5-7}
$n$ & $|x_{n,0}|$ & $|x_{n,1}-1|$ & $|x_{n,2}-2|$ & $|x_{n,0}|$ & $|x_{n,1}-1|$ & $|x_{n,2}-2|$\\
\midrule
$20$ & $9.8\cdot10^{-5}$ & $2.3\cdot10^{-3}$ & $2.1\cdot10^{-2}$ & $5.0\cdot10^{-4}$ & $2.0\cdot10^{-2}$ & $1.8\cdot10^{-1}$\\
$40$ & $2.0\cdot10^{-18}$ & $2.6\cdot10^{-16}$ & $1.8\cdot10^{-14}$ & $1.7\cdot10^{-8}$ & $3.3\cdot10^{-6}$ & $1.7\cdot10^{-4}$\\
$80$ & $6.1\cdot10^{-60}$ & $3.6\cdot10^{-57}$ & $1.2\cdot10^{-54}$ & $4.3\cdot10^{-19}$ & $3.4\cdot10^{-16}$ & $9.1\cdot10^{-14}$\\
$160$ & $1.6\cdot10^{-168}$ & $4.0\cdot10^{-165}$ & $5.7\cdot10^{-162}$ & $7.2\cdot10^{-42}$ & $2.3\cdot10^{-38}$ & $2.7\cdot10^{-35}$\\
$320$ & $4.4\cdot10^{-435}$ & $4.5\cdot10^{-431}$ & $2.7\cdot10^{-427}$ & $8.8\cdot10^{-89}$ & $1.1\cdot10^{-84}$ & $5.7\cdot10^{-81}$\\
\bottomrule
\end{tabular}
\end{table}
